\section{Model Specification: Design Concepts and Parameter Register}
\label{app:specification}

\subsection*{B.1 Design concepts and exclusions}

Table~\ref{tab:design-concepts} gives the model's treatment of the eleven \ac{ODD} design concepts and the eight mechanisms it excludes by design.

\begin{longtable}{L{2.9cm}L{7.3cm}L{3.2cm}}
\caption{\ac{ODD} design concepts and design exclusions}
\label{tab:design-concepts} \\
\multicolumn{3}{p{13.4cm}}{\footnotesize The eleven design concepts of the \ac{ODD} protocol \cite{grimm2020odd}, each with the model's treatment and the literature that grounds it; ``None'' in the source column marks a modelling choice with no external anchor. The final row lists the eight mechanisms excluded by design and the reason for each. The absence of default-to-income and lender-loss feedback limits the contagion mechanisms tested in Section~\ref{sec:results-rq2}.} \\
\toprule
\textbf{Concept} & \textbf{Treatment} & \textbf{Source} \\
\midrule
\endfirsthead
\toprule
\textbf{Concept} & \textbf{Treatment} & \textbf{Source} \\
\midrule
\endhead
\bottomrule
\endfoot
Basic principles & Balance sheets form the accounting core. Heuristic decision rules replace intertemporal optimisation. Statutory regulations act as explicit computational gates. & \cite{cardaci2018inequality, dorazio2017micro, meier2010present, Keys2019, yun2020housing} \\
\addlinespace
Emergence & Default rates and adoption paths follow from household decisions; stacking depths follow from the routing rule of Submodel~13. Borrowing to service existing debt raises future obligations, while lagged peer adoption raises borrowing propensities. Setting $\beta=0$ removes the peer-adoption effect. & None \\
\addlinespace
Adaptation & Households compress discretionary spending and borrow against shortfalls from current liquidity, and their want-driven purchase propensity rises with lagged peer adoption; they form none of the income expectations that D'Orazio and Giulioni give their households. The traditional lender and the platforms are non-adaptive. & \cite{dorazio2017micro} \\
\addlinespace
Objectives & No forward-looking utility. Households satisfy obligations in a fixed priority order, compressing discretionary spending under stress. & \cite{meier2010present} \\
\addlinespace
Learning & Excluded. Behavioural rules are static over the horizon. & None \\
\addlinespace
Prediction & Excluded. Agents treat current income as persistent. & None \\
\addlinespace
Sensing & Agents know their own balance sheets and observe the lagged adoption share among their reference group's \ac{BNPL}-eligible members. In the benchmark the traditional lender sees only bureau-recorded traditional debt and defaults, and each platform sees only its own exposure; the bureau switch adds \ac{BNPL} obligations to the lender's view (Submodel~15). & \cite{nortonrose_bnpl_sa} \\
\addlinespace
Interaction & Households interact with lenders through credit applications, and with each other indirectly through reference-group adoption shares. & \cite{ackert2025bnpl} \\
\pagebreak
Stochasticity & Income shocks and re-employment, archetype and eligibility assignment, want-driven purchases and their size, platform routing order, payment friction, activation order, and threshold endowments (drawn from a separate stream so control arms reproduce exactly across threshold dispersion). Runs are seeded and replicated, with dispersion reported. & None \\
\addlinespace
Collectives & Reference groups (income quintile $\times$ province) aggregate individual adoption into the peer signal. & \cite{cardaci2018inequality} \\
\addlinespace
Observation & Tracked per tick: population default rate; the \ac{CCMR} arrears distribution of credit-active households; traditional debt, arrears and interest; \ac{BNPL} balances, holding, volume and fees; the number of platforms each household owes; reference-group adoption shares; and the share with no liquid savings. Per run: the traditional lender's applications, grants and refusals, order-cap and limit binding rates, and default, purchase size and debt-to-income by income quintile. & None \\
\addlinespace
Exclusions & Eight mechanisms are excluded by design. \emph{No macroeconomic dynamics}: holding the environment fixed isolates the injection effect under baseline conditions. \emph{No contagion}: one aggregate lender with no balance sheet that can fail, and no feedback from household default to income, so default cannot spread through these channels. \emph{No competition among traditional lenders}, which suppresses the lender-conduct mechanism of Hamill et al. \emph{No learning and no scarring}. \emph{No peer effects on consumption}: peer influence acts only on \ac{BNPL} adoption. \emph{Static household composition}. \emph{No credit scores or risk pricing}: the bureau switch acts only through the affordability test. \emph{No other source of credit}: a household that the lender refuses has no informal or alternative lender. & \cite{hamill2023creditcard, cardaci2018inequality} \\
\end{longtable}

\subsection*{B.2 Parameter register}

Table~\ref{tab:param-register} gives every parameter with its value, provenance and sweep range.
It is generated from the parameter declarations in \texttt{simulation/config.py} by
\texttt{notebooks/scripts/generate\_param\_register.py}.

Four groups of entries carry standing flags:
\begin{itemize}
    \item \textbf{The two fitted parameters.} The income-shock probability (Submodel~1) and the
    payment-friction rate (Submodel~6) are fitted to the \ac{CCMR} 90+ and 1--30 day bands, which
    is what makes the baseline calibrated rather than validated
    (Section~\ref{sec:lim-calibration}).
    \item \textbf{The uncited rules.} The shortfall borrowing amount (Submodel~4), \ac{BNPL} relief
    of a shortfall and its cap (Submodel~5), the minimum-payment fraction (Submodel~6), the use of
    credit raised against a committed shortfall (Submodel~7) and platform routing (Submodel~13)
    have no literature anchor. Each is varied in at least one arm. Three assumptions are not
    varied: purchase-size dispersion, the order cap, which binds on 0.3\% of platform requests
    at $\beta=0$, and the payment of \ac{BNPL} instalments before traditional debt
    (Section~\ref{sec:sequence}).
    \item \textbf{The uncalibrated peer coefficients.} $q_{\text{base}}$, $\beta$ and the
    threshold parameters of Submodel~11 have no South African estimate. The
    $\beta = 0$ and $\gamma = 0$ control arms are reported beside every peer-influence result;
    the level-sensitivity arms run at $\beta = 1$ only.
    \item \textbf{The rolling limit.} No South African provider publishes a limit. $\lambda = 0.10$
    is set so that four platform limits sum to about Woolard's illustrative \pounds1{,}000 of
    credit amassed across several lenders \cite{woolard2021}, 39\% of United Kingdom median
    monthly household disposable income, equivalised, in the year to March 2020
    \cite{ons_income_fye2020}. It is swept 0.10--1.0 and reported against a 0.10--0.48 band
    whose upper end expresses Afterpay's maximum of A\$2{,}000 as a share of Australian median
    monthly household disposable income, equivalised, in 2019--20
    \cite{asic2020rep672, abs_income_2019_20}. Equivalised income is below the income of a
    household with more than one member, so both shares overstate the limits relative to the
    household income that $\lambda$ multiplies (Submodel~12).
\end{itemize}

% GENERATED by notebooks/scripts/generate_param_register.py from
% simulation/config.py -- do not edit by hand; re-run the script instead.
\begingroup
\setlength{\tabcolsep}{3pt}
\newcolumntype{R}[1]{>{\scriptsize\raggedright\arraybackslash}p{#1}}
\begin{longtable}{R{2.4cm}R{1.5cm}R{1.75cm}R{1.15cm}R{3.5cm}R{2.85cm}}
\caption{The full parameter register}
\label{tab:param-register} \\
\multicolumn{6}{p{13.4cm}}{\footnotesize Every model parameter. Names, submodels,
provenance and values are read from the parameter declarations in the code, so the
table and the code cannot drift apart. Columns: the parameter name in the code; the
submodel of Table~\ref{tab:submodels} that it belongs to (``design'' marks the
experimental design and ``data'' the data layer); its value in the benchmark
configuration, with both fitted parameters at their fitted values; its provenance; the
source that fixes or motivates the value; and the range swept where one applies.
Provenance: sourced = a citation fixes the value; derived = the value is computed from
data; fitted = the value is chosen by calibration; assumed = no anchor exists. Rand
values are 2017 Rands; the \ac{BNPL} contract parameters are published at 2026 vintage
and deflated once (Appendix~\ref{app:variables}, Part~D.12).} \\
\toprule
\textbf{Parameter} & \textbf{Submodel} & \textbf{Value} & \textbf{Prov.} & \textbf{Source} & \textbf{Sweep} \\
\midrule
\endfirsthead
\toprule
\textbf{Parameter} & \textbf{Submodel} & \textbf{Value} & \textbf{Prov.} & \textbf{Source} & \textbf{Sweep} \\
\midrule
\endhead
\bottomrule
\endfoot
\texttt{n\_\allowbreak{}ticks} & design & 52 & sourced & A 24-month horizon at a 14-day tick (Section~\ref{sec:sim-params}). & -- \\
\addlinespace
\texttt{burn\_\allowbreak{}in} & design & 12 & sourced & The first 12 ticks are excluded from post-burn-in summaries (Section~\ref{sec:sim-params}). & -- \\
\addlinespace
\texttt{n\_\allowbreak{}agents} & data & 5{,}000 & derived & The full resampled population of 5{,}000 (Appendix~\ref{app:variables}, Part~D.10). & 1{,}000 / 5{,}000 / 10{,}000 \\
\addlinespace
\texttt{shock\_\allowbreak{}prob} & 1 & 0.016 & fitted & Fitted to the \ac{CCMR} 90+ day band. Reported against the \ac{QLFS} 2017 job-separation band of 0.54\%--1.17\% per tick, to which it is not fitted. & calibration grid; 0.005--0.026 in the Sobol design \\
\addlinespace
\texttt{shock\_\allowbreak{}persistent} & 1 & on & sourced & An unemployment spell lasts until re-employment. In the \ac{QLFS} 2017, 68.4\% of the unemployed were still unemployed a quarter later. & single-tick shock as a robustness arm \\
\addlinespace
\texttt{shock\_\allowbreak{}exit\_\allowbreak{}prob} & 1 & 0.018724 & sourced & Re-employment hazard per tick. In the \ac{QLFS} 2017, 11.6\% of the unemployed moved into employment in a quarter, which is 1.87\% per tick. & -- \\
\addlinespace
\texttt{discretionary\_\allowbreak{}floor} & 2 & 0 & sourced & Discretionary spending can be compressed to zero. & -- \\
\addlinespace
\texttt{q\_\allowbreak{}base} & 3/11 & 0.05 & assumed & Probability of a want-driven purchase in a tick without peer influence. No source. & 0.01--0.30 in the Sobol design \\
\addlinespace
\texttt{beta} & 11 & 0 & assumed & Strength of peer influence under the linear rule (Equation~\ref{eq:peer}). No source. $\beta = 0$ is the control. & 0 / 0.5 / 1 / 2 / 3 in the platform and access experiments; 0 and 1 elsewhere \\
\addlinespace
\texttt{peer\_\allowbreak{}mechanism} & 11 & linear & sourced & The rule by which adoption spreads: linear (Equation~\ref{eq:peer}) or heterogeneous thresholds \cite{granovetter1978threshold}. & both rules on the same access grid \\
\addlinespace
\texttt{mu\_\allowbreak{}theta} & 11 & 0.3 & assumed & Mean adoption threshold: the share of a household's reference group that must hold \ac{BNPL} before the household responds. No source. & 0.15 / 0.30 / 0.45 at full access \\
\addlinespace
\texttt{sigma\_\allowbreak{}theta} & 11 & 0.2 & assumed & Dispersion of adoption thresholds, truncated to $[0,1]$. No source. Granovetter's argument concerns the dispersion of thresholds, so this is the axis swept. & 0.05--0.40 \\
\addlinespace
\texttt{gamma} & 11 & 0 & assumed & Rise in the purchase probability once a household's threshold is crossed. No source. $\gamma = 0$ is the control. & 0 and 0.3 \\
\addlinespace
\texttt{amount\_\allowbreak{}rule} & 4 & shortfall & assumed & No estimate was located of what households borrow against a shortfall. & exact shortfall / $+25\%$ / plus one tick of committed expenditure \\
\addlinespace
\texttt{shortfall\_\allowbreak{}bnpl\_\allowbreak{}capped} & 5 & on & assumed & Limits the \ac{BNPL} share of a shortfall request to $\kappa$ times the household's monthly discretionary budget, the quantity that sizes a want-driven purchase. & capped / uncapped, crossed with the three amount rules \\
\addlinespace
\texttt{bnpl\_\allowbreak{}purchase\_\allowbreak{}base} & 4/12 & discretionary & sourced & The budget against which a purchase is sized. \ac{BNPL} finances discretionary consumption, so the reference is the household's discretionary budget. & discretionary / income \\
\addlinespace
\texttt{bnpl\_\allowbreak{}purchase\_\allowbreak{}ratio} & 4/12 & 0.1387 & derived & $\kappa$: the share of discretionary spending in the categories that \ac{BNPL} finances, derived from \ac{IES} 2022/23 microdata \cite{statssa_ies_2023}. & 0.07--0.28, half to twice the derived share \\
\addlinespace
\texttt{bnpl\_\allowbreak{}purchase\_\allowbreak{}cv} & 4 & 0.6 & assumed & Dispersion of purchase size around its mean, which is lognormal. No source. & -- \\
\addlinespace
\texttt{min\_\allowbreak{}payer\_\allowbreak{}share} & 6 & 0.29 & sourced & 29\% of United States card accounts pay at or near the minimum \cite{Keys2019}. & 0.20--0.40 \\
\addlinespace
\texttt{payment\_\allowbreak{}friction} & 6 & 0.09 & fitted & Probability per tick of missing a traditional instalment that the household could pay \cite{Kuchler2021}. Fitted to the \ac{CCMR} 1--30 day band. & calibration grid \\
\addlinespace
\texttt{min\_\allowbreak{}payment\_\allowbreak{}frac} & 6 & 0.05 & assumed & Minimum payment as a share of the balance per month. The \ac{NCA} prescribes no formula. & 0.025--0.10 \\
\addlinespace
\texttt{k\_\allowbreak{}default} & 7 & 7 & sourced & Seven ticks are 98 days, the first tick boundary past the 90-day impairment convention. & four ticks (56 days) \\
\addlinespace
\texttt{bnpl\_\allowbreak{}bureau\_\allowbreak{}visible} & 10/15 & off & sourced & \ac{BNPL} is not reported to the bureau in the benchmark \cite{nortonrose_bnpl_sa, transunion_cps_sa_2025}. & on / off (Section~\ref{sec:scenarios}) \\
\addlinespace
\texttt{bnpl\_\allowbreak{}enabled} & design & off & sourced & The 2017 baseline has no \ac{BNPL}; the experiments enable it. & on / off \\
\addlinespace
\texttt{bnpl\_\allowbreak{}access\_\allowbreak{}rate} & 12 & 1 & derived & Share of banked households that are eligible for \ac{BNPL}. Banked households are 82.8\% of the population. & 0--1.0 in seven steps \\
\addlinespace
\texttt{n\_\allowbreak{}platforms} & 13 & 4 & sourced & PayJustNow, Payflex, Mobicred and TymeBank were active in South Africa. & 1--6 \\
\addlinespace
\texttt{bnpl\_\allowbreak{}order\_\allowbreak{}cap} & 12 & 9{,}926.83 & assumed & A per-order cap of R15{,}000 at 2026 prices, deflated to 2017 Rands. Payflex's terms state no maximum order value \cite{payflex_terms}. & -- \\
\addlinespace
\texttt{bnpl\_\allowbreak{}limit\_\allowbreak{}income\_\allowbreak{}multiple} & 12 & 0.1 & assumed & $\lambda$: the limit per platform as a multiple of monthly income. No South African provider publishes a limit; the value is set against Woolard's illustrative stacked exposure \cite{woolard2021, ons_income_fye2020}. & 0.10--1.0 \\
\addlinespace
\texttt{bnpl\_\allowbreak{}instalments} & 14 & 4 & sourced & Pay-in-four: a quarter at checkout and a quarter at each of the next three ticks \cite{payflex_terms}. & -- \\
\addlinespace
\texttt{bnpl\_\allowbreak{}late\_\allowbreak{}fee\_\allowbreak{}per\_\allowbreak{}tick} & 14 & 125.74 & sourced & A late fee of R95 a week, R190 a tick, at 2026 vintage \cite{payflex_terms}, deflated to 2017 Rands. & -- \\
\addlinespace
\texttt{bnpl\_\allowbreak{}late\_\allowbreak{}fee\_\allowbreak{}cap} & 14 & 188.61 & sourced & Late fees on an agreement are capped at R285 over its life, at 2026 vintage \cite{payflex_terms}, deflated to 2017 Rands. & -- \\
\addlinespace
\texttt{bnpl\_\allowbreak{}late\_\allowbreak{}fee\_\allowbreak{}cap\_\allowbreak{}share} & 14 & 0.5 & sourced & Late fees on an agreement are capped at the lower of the Rand cap and half the purchase price \cite{payflex_terms}. & -- \\
\addlinespace
\texttt{bnpl\_\allowbreak{}affordability\_\allowbreak{}check} & 15 & off & sourced & Applies the Regulation 23A test to each \ac{BNPL} request, the analogue of the \ac{FCA}'s affordability check \cite{fca_ps26_1}. & on / off (Section~\ref{sec:scenarios}) \\
\addlinespace
\texttt{k\_\allowbreak{}cool} & 15 & 0 & sourced & Ticks for which want-driven purchases are blocked after a want-driven purchase. One tick is the fourteen-day right of withdrawal in the United Kingdom's rules \cite{fca_ps26_1}. & 0--4 ticks (Appendix~\ref{app:levers}) \\
\addlinespace
\texttt{stacking\_\allowbreak{}cap} & 15 & none & assumed & Cap on platforms owed: a household that owes this many platforms or more is refused every new agreement. Hypothetical; no jurisdiction imposes one. & 1--3 (Appendix~\ref{app:levers}) \\
\addlinespace
\texttt{activation} & 17 & random & sourced & Random order, redrawn every tick \cite{comer2013activation, alizadeh2015activation}. & fixed order as a robustness arm \\
\end{longtable}
\endgroup

